%% appendix_geo_scores.tex — Geomorphological Organic Abundance Scores
%% Extracted from main.tex
\chapter{Geomorphological Organic Abundance Scores}
\label{app:geo_scores}

\begin{table}[H]
\centering
\begin{tabular}{llp{5.0cm}c}
\toprule
\textbf{Class} & \textbf{Name} & \textbf{Physical basis} & \textbf{Score} \\
\midrule
Dunes     & Equatorial sand seas & Tholin-dominated; highest VIMS ratio \citep{Lorenz2006}           & 0.82 \\
Plains    & Smooth lowlands      & Organic sediment blanket \citep{Lopes2019b}                        & 0.68 \\
Basins    & Topographic lows     & Organic accumulation \citep{Lopes2020}                            & 0.58 \\
Labyrinth & Dissected terrain    & Spectral affinity with plains \citep{Malaska2020}                 & 0.55 \\
Craters   & Impact structures    & Mixed water-ice and organic ejecta \citep{Neish2018}              & 0.35 \\
Mountains & Ridge terrain        & Water-ice-rich crustal material \citep{Solomonidou2018}           & 0.25 \\
Lakes     & Hydrocarbon seas     & Liquid hydrocarbons; low solid organic content \citep{Brown2008}  & 0.05 \\
\bottomrule
\end{tabular}
\caption{Organic abundance proxy scores used for the geomorphological
blended VIMS mode of Feature~2.  Grounded in published VIMS spectral
characterisation of each terrain class.  This lookup includes a Lakes row,
which feeds the liquid-hydrocarbon ($f_1$) and organic-abundance ($f_2$)
features; the lakes class is absent from the six-class geomorphology raster
that underlies the geomorphological-diversity feature ($f_7$,
Section~\ref{sec:feat_geodiv}), so the seven rows here and the six terrain
classes used by $f_7$ are consistent rather than contradictory.}
\label{tab:geo_scores}
\end{table}